% ===================================================================
% data/schaetzung_teil5.tex
% Eigene Schaetzung des Verfassers zu Investitionsbedarf, Kapitalbindung
% und Liquiditaetsbedarf (jaehrlich) der drei Handlungsoptionen in
% Abschnitt 5 - UNABHAENGIG von den fuer Abschnitt 6 vorgegebenen
% Betraegen in data/cf_makros.tex. Die Aufgabenstellung erlaubt und
% erwartet diese Abweichung ausdruecklich ("ggf. abweichend von Ihren
% vorausgegangenen Analysen", Aufgabe 6).
%
% Methode, Anker und Beleg: docs/prompt-teil5-eigene-schaetzung.md.
% Von Hand hergeleitet, nicht generiert - anders als data/cf_makros.tex.
% ===================================================================

% --- Kernrate: aktivierbare Entwicklungskosten ----------------------
% \FEQuoteZFIII (1,5 %, GB 2023 S. 14, ertragsstarkes Jahr) x
% \UmsatzZFIV (312,346 Mio., Konzernumsatz 2024) x 3 Jahre = 14,06,
% gerundet 14,1 Mio. Deckt nur den aktivierungsfaehigen
% Personalkostenanteil der Entwicklung ab (Aktivierungskriterien nach
% IAS 38), nicht Vertrieb, Anwendungstechnik oder Zertifizierung eines
% Markteintritts - daher als Kern behandelt und je Option mit einem
% Markterschliessungsfaktor skaliert, nicht direkt uebernommen.
% (0.015 * 312.346 * 3 = 14.05557)

% --- Option A: Diversifikation in drei branchenfremde Zielmaerkte --
% (Clean Tech, Aerospace, Life Sciences ueber das Segment Next
% Automation). Faktor 3 auf die Kernrate: kein bestehender
% Kundenzugang in diesen Branchen, daher voller Vertriebs- und
% Zertifizierungsaufwand zusaetzlich zur reinen Entwicklung.
\newcommand{\InvestitionsbedarfEigenA}{42.2}  % 14.05557 * 3
\newcommand{\KapitalbindungEigenA}{10.6}      % InvestitionsbedarfEigenA * 25 % (Umlaufvermoegen/Projektvolumen bei Auftragsfertigern)
\newcommand{\LiquiditaetsbedarfEigenA}{14.1}  % InvestitionsbedarfEigenA / 3 Jahre, gleichmaessig verteilt

% --- Option B: Erweiterung entlang der bestehenden Wertschoepfung --
% (Batterie-/Brennstoffzellenfertigung im bestehenden E-Mobility-
% Umfeld). Faktor 1,5 auf die Kernrate statt 3: Verfahren, Kunden und
% Vertrieb sind bereits bekannt, ein Teil der Kompetenz ist mit der
% Aumann Lauchheim GmbH bereits im Haus.
\newcommand{\InvestitionsbedarfEigenB}{21.1}  % 14.05557 * 1.5
\newcommand{\KapitalbindungEigenB}{5.3}       % InvestitionsbedarfEigenB * 25 %
\newcommand{\LiquiditaetsbedarfEigenB}{7.0}   % InvestitionsbedarfEigenB / 3 Jahre

% --- Option C: Akquisition ------------------------------------------
% Kein F&E-Anker moeglich: das Zielobjekt ist laut Aufgabenstellung
% nicht bestimmt, daher keine Entwicklungsquote anwendbar. Stattdessen
% eine Kapazitaetsgrenze: ein einzelner, weitgehend unumkehrbarer
% Kapitaleinsatz sollte die Nettofinanzliquiditaet 2024
% (\NettoliquiditaetZFIV, 138,2 Mio.) nicht zur Haelfte uebersteigen,
% um Handlungsfaehigkeit fuer das laufende Geschaeft und weitere
% Gelegenheiten zu erhalten (0.5 * 138.2 = 69.1 Mio. Gesamtobergrenze).
% Aufgeteilt im selben Verhaeltnis wie bei A und B (Kapitalbindung =
% 25 % des Investitionsbedarfs, also Investitionsbedarf = Gesamt / 1.25).
\newcommand{\InvestitionsbedarfEigenC}{55.3}  % 69.1 / 1.25
\newcommand{\KapitalbindungEigenC}{13.8}      % 69.1 - 55.3 (= InvestitionsbedarfEigenC * 25 %)
\newcommand{\LiquiditaetsbedarfEigenC}{69.1}  % InvestitionsbedarfEigenC + KapitalbindungEigenC, im Vollzugsjahr faellig, nicht verteilt (Kaufpreis wird nicht in Raten gezahlt)
